% GENERATED FILE. Edit canonical lessons, then run npm run generate:books.
% canonical-source-hash: fnv1a64:e2aad2acca751200
% canonical-lessons: GU-C367-sambhavatah, GU-C367-acuk, GU-C367-niscit, GU-C367-dharna, GU-C367-akvu

\chapter{Probably}
\label{ch:gu-probably}
\hlchaptermodality{367}

\begin{chapteropening}
\textbf{By the end of this chapter:} \emph{I can say how likely something is.}

Complete the last lesson of chapter 367: I can say how likely something is.
\end{chapteropening}

\section[\texorpdfstring{\gu{સંભવતઃ} (sambhavataḥ)}{સંભવતઃ (sambhavataḥ)}]{\gu{સંભવતઃ} (sambhavataḥ) — probably}

\label{lesson:GU-C367-sambhavatah}



\begin{quote}
Before the new one: say the Gujarati for a conjecture, then the Gujarati for to want.
\end{quote}

\subsection*{You'll want to know: \gu{સંભવતઃ}}
\textbf{\gu{સંભવતઃ}} — \emph{sambhavataḥ} — \textquotedblleft{}probably\textquotedblright{}.

\begin{quote}
\gu{જો} \gu{ટ્રેન} \gu{સમયસર} \gu{હશે} \gu{તો} \gu{અમે} \gu{સંભવતઃ} \gu{સાંજે} \gu{પહોંચીશું}.
\end{quote}

\begin{quote}
\emph{jo ṭren samaysar haśe to ame sambhavataḥ sā̃je pahõcīśũ.}
\end{quote}

\begin{quote}
If the train is on time, we will probably arrive by evening.
\end{quote}

\begin{cousinweb}
\textbf{\gu{સંભવતઃ}} (\emph{sambhavataḥ}) is built from Sanskrit \emph{sambhava}, \textquotedblleft{}coming about, possibility\textquotedblright{} (Monier-Williams), and the adverb ending \emph{-taḥ}: by way of what is possible.
\end{cousinweb}

\begin{grammarlens}[title={What you've built}]
One more word for saying what happens if.
\end{grammarlens}

\subsection*{Your turn}
\begin{itemize}
\raggedright
  \item \emph{Say it:} \emph{sambhavataḥ}
  \item \emph{Say it:} \emph{sambhavataḥ}, once more
  \item \emph{Say it:} \emph{jo ṭren samaysar haśe to ame sambhavataḥ sā̃je pahõcīśũ.} — If the train is on time, we will probably arrive by evening.
  \item \emph{Recall:} say the Gujarati for a conjecture, then the Gujarati for to want, then say \emph{sambhavataḥ} again
\end{itemize}

\begin{tcolorbox}[breakable,colback=teal!4,colframe=teal!35!black,title={Before you move on}]
What does \gu{સંભવતઃ} mean? (\textquotedblleft{}Probably\textquotedblright{}.) Say it once more.
\end{tcolorbox}
\section[\texorpdfstring{\gu{અચૂક} (acūk)}{અચૂક (acūk)}]{\gu{અચૂક} (acūk) — without fail}

\label{lesson:GU-C367-acuk}



\begin{quote}
Before the new one: say the Gujarati for to want, then the Gujarati for probably.
\end{quote}

\subsection*{You'll want to know: \gu{અચૂક}}
\textbf{\gu{અચૂક}} — \emph{acūk} — \textquotedblleft{}without fail\textquotedblright{}.

\begin{quote}
\gu{જો} \gu{તમે} \gu{આ} \gu{દવા} \gu{લેશો} \gu{તો} \gu{તાવ} \gu{અચૂક} \gu{મટશે}.
\end{quote}

\begin{quote}
\emph{jo tame ā davā leśo to tāv acūk maṭśe.}
\end{quote}

\begin{quote}
If you take this medicine, the fever will go without fail.
\end{quote}

\begin{cousinweb}
\textbf{\gu{અચૂક}} (\emph{acūk}) is built from \emph{a-}, \textquotedblleft{}not\textquotedblright{}, and \emph{cūk}, \textquotedblleft{}a slip, a miss\textquotedblright{}, from a reconstructed \emph{cukk}, \textquotedblleft{}to fall short\textquotedblright{} (CDIAL 4848): with no miss.
\end{cousinweb}

\begin{grammarlens}[title={What you've built}]
One more word for saying what happens if.
\end{grammarlens}

\subsection*{Your turn}
\begin{itemize}
\raggedright
  \item \emph{Say it:} \emph{acūk}
  \item \emph{Say it:} \emph{acūk}, once more
  \item \emph{Say it:} \emph{jo tame ā davā leśo to tāv acūk maṭśe.} — If you take this medicine, the fever will go without fail.
  \item \emph{Recall:} say the Gujarati for to want, then the Gujarati for probably, then say \emph{acūk} again
\end{itemize}

\begin{tcolorbox}[breakable,colback=teal!4,colframe=teal!35!black,title={Before you move on}]
What does \gu{અચૂક} mean? (\textquotedblleft{}Without fail\textquotedblright{}.) Say it once more.
\end{tcolorbox}
\section[\texorpdfstring{\gu{નિશ્ચિત} (niścit)}{નિશ્ચિત (niścit)}]{\gu{નિશ્ચિત} (niścit) — definite; settled}

\label{lesson:GU-C367-niscit}



\begin{quote}
Before the new one: say the Gujarati for probably, then the Gujarati for without fail.
\end{quote}

\subsection*{You'll want to know: \gu{નિશ્ચિત}}
\textbf{\gu{નિશ્ચિત}} — \emph{niścit} — \textquotedblleft{}definite; settled\textquotedblright{}.

\begin{quote}
\gu{જો} \gu{તારીખ} \gu{નિશ્ચિત} \gu{થશે} \gu{તો} \gu{હું} \gu{તમને} \gu{ફોન} \gu{કરીશ}.
\end{quote}

\begin{quote}
\emph{jo tārīkh niścit thaśe to hũ tamne phon karīś.}
\end{quote}

\begin{quote}
If the date is settled, I will phone you.
\end{quote}

\begin{cousinweb}
\textbf{\gu{નિશ્ચિત}} (\emph{niścit}) is a tatsama, borrowed from Sanskrit \emph{niścita}, \textquotedblleft{}determined, settled\textquotedblright{} (Monier-Williams).
\end{cousinweb}

\begin{grammarlens}[title={What you've built}]
One more word for saying what happens if.
\end{grammarlens}

\subsection*{Your turn}
\begin{itemize}
\raggedright
  \item \emph{Say it:} \emph{niścit}
  \item \emph{Say it:} \emph{niścit}, once more
  \item \emph{Say it:} \emph{jo tārīkh niścit thaśe to hũ tamne phon karīś.} — If the date is settled, I will phone you.
  \item \emph{Recall:} say the Gujarati for probably, then the Gujarati for without fail, then say \emph{niścit} again
\end{itemize}

\begin{tcolorbox}[breakable,colback=teal!4,colframe=teal!35!black,title={Before you move on}]
What does \gu{નિશ્ચિત} mean? (\textquotedblleft{}Definite; settled\textquotedblright{}.) Say it once more.
\end{tcolorbox}
\section[\texorpdfstring{\gu{ધારણા} (dhārṇā)}{ધારણા (dhārṇā)}]{\gu{ધારણા} (dhārṇā) — an assumption}

\label{lesson:GU-C367-dharna}



\begin{quote}
Before the new one: say the Gujarati for without fail, then the Gujarati for definite.
\end{quote}

\subsection*{You'll want to know: \gu{ધારણા}}
\textbf{\gu{ધારણા}} — \emph{dhārṇā} — \textquotedblleft{}an assumption\textquotedblright{}.

\begin{quote}
\gu{જો} \gu{મારી} \gu{ધારણા} \gu{ખોટી} \gu{હશે} \gu{તો} \gu{હું} \gu{માફી} \gu{માગીશ}.
\end{quote}

\begin{quote}
\emph{jo mārī dhārṇā khoṭī haśe to hũ māphī māgīś.}
\end{quote}

\begin{quote}
If my assumption is wrong, I will apologise.
\end{quote}

A feminine noun.

\begin{cousinweb}
\textbf{\gu{ધારણા}} (\emph{dhārṇā}) is a tatsama, borrowed from Sanskrit \emph{dhāraṇā}, \textquotedblleft{}holding, bearing, keeping in mind\textquotedblright{} (Monier-Williams).
\end{cousinweb}

\begin{grammarlens}[title={What you've built}]
One more word for saying what happens if.
\end{grammarlens}

\subsection*{Your turn}
\begin{itemize}
\raggedright
  \item \emph{Say it:} \emph{dhārṇā}
  \item \emph{Say it:} \emph{dhārṇā}, once more
  \item \emph{Say it:} \emph{jo mārī dhārṇā khoṭī haśe to hũ māphī māgīś.} — If my assumption is wrong, I will apologise.
  \item \emph{Recall:} say the Gujarati for without fail, then the Gujarati for definite, then say \emph{dhārṇā} again
\end{itemize}

\begin{tcolorbox}[breakable,colback=teal!4,colframe=teal!35!black,title={Before you move on}]
What does \gu{ધારણા} mean? (\textquotedblleft{}An assumption\textquotedblright{}.) Say it once more.
\end{tcolorbox}
\section[\texorpdfstring{\gu{આંકવું} (ā̃kvũ)}{આંકવું (ā̃kvũ)}]{\gu{આંકવું} (ā̃kvũ) — to assess}

\label{lesson:GU-C367-akvu}



\begin{quote}
Before the new one: say the Gujarati for definite, then the Gujarati for an assumption.
\end{quote}

\subsection*{You'll want to know: \gu{આંકવું}}
\textbf{\gu{આંકવું}} — \emph{ā̃kvũ} — \textquotedblleft{}to assess\textquotedblright{}.

\begin{quote}
\gu{જો} \gu{તમે} \gu{ઘરની} \gu{કિંમત} \gu{ઓછી} \gu{આંકશો} \gu{તો} \gu{ખોટ} \gu{જશે}.
\end{quote}

\begin{quote}
\emph{jo tame gharnī kimmat ochī ā̃kśo to khoṭ jaśe.}
\end{quote}

\begin{quote}
If you value the house too low, you will make a loss.
\end{quote}

Past, with \gu{તેણે} (he or she …): \gu{તેણે} \gu{ઘરની} \gu{કિંમત} \gu{આંકી}, \emph{teṇe gharnī kimmat ā̃kī}.

\begin{cousinweb}
\textbf{\gu{આંકવું}} (\emph{ā̃kvũ}) is inherited from Sanskrit \emph{aṅkáyati}, \textquotedblleft{}marks\textquotedblright{}; CDIAL lists the Gujarati verb, \textquotedblleft{}to mark, to appraise\textquotedblright{} (CDIAL 104).
\end{cousinweb}

\begin{grammarlens}[title={What you've built}]
One more word for saying what happens if.
\end{grammarlens}

\subsection*{Your turn}
\begin{itemize}
\raggedright
  \item \emph{Say it:} \emph{ā̃kvũ}
  \item \emph{Say it:} \emph{ā̃kvũ}, once more
  \item \emph{Say it:} \emph{jo tame gharnī kimmat ochī ā̃kśo to khoṭ jaśe.} — If you value the house too low, you will make a loss.
  \item \emph{Recall:} say the Gujarati for definite, then the Gujarati for an assumption, then say \emph{ā̃kvũ} again
\end{itemize}

\begin{tcolorbox}[breakable,colback=teal!4,colframe=teal!35!black,title={Before you move on}]
What does \gu{આંકવું} mean? (\textquotedblleft{}To assess\textquotedblright{}.) Say it once more.
\end{tcolorbox}
